\documentclass[11pt]{article}
\usepackage{amsmath,amssymb}
\usepackage[margin=1in]{geometry}
\usepackage[hidelinks]{hyperref}
\title{Research summary: 4--7 October 2026}
\author{Daniel}
\date{7 October 2026}
\begin{document}
\maketitle

\section*{Gr\"unbaum: understanding the v2 argument}

The main objective was to understand the
argument for deforming the Gr\"unbaum
Stanley--Reisner scheme. The
manuscripts were organized into v0, v1 and v2.
The new AI-written v2 draft treats equivariant
formal lifting over
$$
A=\mathbb{Z}_{(101)}.
$$
Its stated result is that
$$
(T^2_0)^G=0,\qquad G\cong S_3,
$$
and that a prescribed equivariant flat graded
deformation of finite order extends to all
orders. This is a draft to be understood and
checked; formal lifting alone does not
establish smoothness of a general fibre.

My own notes isolate the mechanism behind the
invariant-vanishing argument. The ordinary
degree-zero obstruction module decomposes into
multigraded pieces, with nine orbits of three
nonzero pieces in the draft calculation.
Each piece has the form
$$
A^2/A(1,1).
$$
Swapping the two components acts by minus one.
For the piece labelled $a^2/(cg)$, the
stabilizing involution
$$
\sigma=(b\,h)(c\,g)(d\,f)
$$
exchanges the two components. An invariant
class would therefore satisfy $x=-x$, hence
$x=0$ because $2$ is invertible. The full
argument requires the corresponding check
for every orbit of pieces.

The personal notes still need one coefficient
correction: $A$ in the v2 draft is the
localization $\mathbb{Z}_{(101)}$, not the
field with 101 elements. Reduction modulo
101 gives $\mathbb{F}_{101}$.

A separate linkage notebook was created and
research correspondence archived. The linkage
notebook currently contains the problem
statement and an empty notes section.

\section*{Differentials and the cotangent complex}

In HP, I rebuilt the route from K\"ahler
differentials to the cotangent complex.
For a polynomial presentation $R=P/I$, the
starting point is the conormal sequence
$$
I/I^2\longrightarrow
\Omega^1_{P/k}\otimes_P R
\longrightarrow\Omega^1_{R/k}
\longrightarrow0.
$$
The discussion explains base change to $R$,
why the differentials of relations vanish,
and why taking the cokernel gives the
differentials of the quotient algebra.

The cotangent-complex introduction was
reworked around a simplicial polynomial
resolution of $R$. Applying differentials,
base-changing to $R$, and passing to the
associated complex gives $L_{R/k}$. This is
the background needed for
$$
T^i_{R/k}
=\operatorname{Ext}^i_R(L_{R/k},R).
$$

\section*{First-order algebra deformations}

I worked through a first-order deformation
as a flat algebra $\widetilde R$ over
$$
D=k[\varepsilon]/(\varepsilon^2),
\qquad
\widetilde R\otimes_D k\cong R.
$$
Flatness identifies
$\varepsilon\widetilde R$ with $R$.
The kernel is square-zero, giving
$$
0\longrightarrow\varepsilon\widetilde R
\longrightarrow\widetilde R
\longrightarrow R\longrightarrow0.
$$
This sequence splits as $k$-vector spaces.
It is not canonically a sequence of
$R$-modules: an algebra map
$R\to\widetilde R$ is not given.

After choosing a $k$-linear section, the
difference between $s(a)s(b)$ and $s(ab)$
defines a multiplication correction
$\mu(a,b)\in R$. The exercises derive
bilinearity, symmetry and the associativity
condition
$$
a\mu(b,c)-\mu(ab,c)
+\mu(a,bc)-\mu(a,b)c=0.
$$
Changing the section by a linear map $h$
changes the correction by
$$
\mu'(a,b)-\mu(a,b)
=a h(b)-h(ab)+h(a)b.
$$
Thus the calculations give cocycles modulo
coboundaries. They were connected to
Hochschild cohomology for associative
deformations, while retaining the symmetry
condition for commutative deformations.

The current notes begin the bridge to
cotangent cohomology: a polynomial algebra
$P_0$ mapping onto $R$ lifts to a square-zero
extension by lifting its generators.
The full proof that algebra extensions by
$M$ are classified by
$\operatorname{Ext}^1_R(L_{R/k},M)$ remains
the next step.

\section*{Module extensions and homology}

I rewrote the module-extension discussion
around the derived-functor construction of
$\operatorname{Ext}^1_R(B,A)$ and its
classification of extensions
$$
0\longrightarrow A\longrightarrow E
\longrightarrow B\longrightarrow0.
$$
The revisions cover lifting from a projective
resolution, equivalence of extensions,
the zero class, and Baer addition.
Splitting exercises explain the equivalent
section/retraction criteria, why vector-space
extensions split, and why
$\operatorname{Ext}^1_R(R,R)=0$.

I also reconstructed the missed homology
lectures: simplicial approximation and
contiguity, acyclic carriers, singular
homology and homotopy invariance, comparison
with simplicial homology, coefficients,
and the universal coefficient theorem.

\section*{Next mathematical step}

Complete the passage from square-zero
algebra extensions to cotangent cohomology,
then return to the multigraded obstruction
calculation and the equivariant lifting
argument in v2. Understanding and checking
that draft, and establishing smoothness,
remain separate tasks.

\end{document}
